\sect{Свёрточный модуль и обработка граничных случаев}{cnn-corner}
%
В аудиорежиме гауссова эмиссионная модель HMM становится неадекватной:
аккорд из нескольких одновременно звучащих нот, артефакты педали и
тембральная окраска инструмента дают многомодальные распределения
наблюдений, плохо приближаемые единственным гауссовским законом. Для
оценки апостериорного распределения по высотам нот мы используем
лёгкую свёрточную сеть, выполняющую роль обучаемой эмиссионной модели.

Архитектура сети представлена на \Figref{cnn-arch}. Входом служит
лог-мел-спектрограмма размером $128\times 40$, рассчитанная по окну
$23{,}2$\,мс с шагом $11{,}6$\,мс (что соответствует одному
блоку синтеза, см.~\Secref{realtime}). За четырьмя свёрточными
блоками с фильтрами $3\times 3$, числом каналов $32$, $64$, $128$,
$256$ и слоями ReLU и max-pooling следует глобальный усредняющий
слой и полносвязный выход размерности~$88$ (число клавиш фортепиано).
Полное число параметров составляет $0{,}9$\,млн, что обеспечивает
инференс за $1{,}8$\,мс на CPU Intel~i7 десятого поколения.

\par\addvspace{2pt}
{\centering
\refstepcounter{figure}\figlab{cnn-arch}%
\resizebox{\columnwidth}{!}{% =====================================================================
% TikZ figure: CNN front-end architecture (BLACK & WHITE only)
% Compact 4+4 layout to fit a single REVTeX two-column column.
% Used in: sections/06-cnn-and-corner-cases.tex
% =====================================================================
\begin{tikzpicture}[
    node distance = 3mm and 2.5mm,
    layer/.style={rectangle, draw=black, thick,
                  minimum height=10mm, minimum width=10mm,
                  font=\tiny, align=center},
    arrow/.style={-Stealth, thick}]

  \node[layer]                 (in)   {Mel\\$128{\times}40$};
  \node[layer, right=of in,
        pattern=north east lines, pattern color=black!30]
                               (c1)   {Conv\\$3{\times}3$\\32};
  \node[layer, right=of c1,
        pattern=north east lines, pattern color=black!30]
                               (c2)   {Conv\\$3{\times}3$\\64};
  \node[layer, right=of c2,
        pattern=north east lines, pattern color=black!30]
                               (c3)   {Conv\\$3{\times}3$\\128};

  \node[layer, below=7mm of c3,
        pattern=north east lines, pattern color=black!30]
                               (c4)   {Conv\\$3{\times}3$\\256};
  \node[layer, left=of c4]     (gap)  {GAP};
  \node[layer, left=of gap]    (fc)   {FC\\$\to{}$128};
  \node[layer, left=of fc]     (sm)   {Soft-\\max\\88};

  \foreach \a/\b in {in/c1, c1/c2, c2/c3} \draw[arrow] (\a) -- (\b);
  \draw[arrow] (c3) -- (c4);
  \foreach \a/\b in {c4/gap, gap/fc, fc/sm} \draw[arrow] (\a) -- (\b);

\end{tikzpicture}
}\par
\vspace{2pt}
\footnotesize
\figurename~\thefigure. Архитектура свёрточного фронтенда.
Заштрихованы свёрточные блоки, GAP~--- глобальный усредняющий слой,
FC~--- полносвязный слой, softmax~--- выход по клавишам фортепиано.\par}
\addvspace{4pt}

Обучение проводилось на смеси MAPS и собственного датасета
\textit{CU-Concerto-2026} с лоссом
\begin{equation}
\mathcal{L}_{\text{CNN}} = -\frac{1}{T}\sum_{t=1}^{T}
\sum_{k=1}^{88} y_{t,k}\,\log P_{\text{CNN}}(k\mid o_t),
\eqlab{cnn-loss}
\end{equation}
где $y_{t,k}\in\{0,1\}$ обозначает наличие $k$-й ноты в момент~$t$.
Оптимизатор Adam с шагом $3\!\times\!10^{-4}$, batch-size~$64$,
обучение продолжалось до сходимости валидационной F1-метрики
($24$ эпохи). Полученное распределение $P_{\text{CNN}}(\cdot\mid o_t)$
подставляется в~\Eqref{forward-recursion} вместо гауссовой эмиссии:
$b_i(o_t)\equiv P_{\text{CNN}}(p_i\mid o_t)$. Это снижает ошибку
выравнивания на полифонических фрагментах в~$1{,}6$ раза по
сравнению с базовой гауссовой моделью.

Граничные случаи отдельно классифицированы и обрабатываются явными
правилами. Долгие ферматы и педальные удержания проявляются как
большие значения $\alpha_t(\hat\varphi)$ при отсутствии новых
сильных onset-кандидатов; в этом случае счётчик длительности
$d$ растёт без перехода на следующее состояние, а темп оркестра
плавно замедляется до $0{,}7$ номинального. Резкие ритардандо в
каденции выявляются по знаковому изменению производной темпа в
окне $250$\,мс и инициируют включение специального профиля растяжения
с увеличенной чувствительностью. Полные пропуски тактов
обрабатываются процедурой повторной синхронизации, описанной
в~\Secref{realtime}: HMM переинициализируется на участке
$\pm 4$ ноты от наилучшей текущей оценки.

Орнаменты (форшлаги, мелизмы) представляют отдельную сложность:
их быстрые ноты не предусмотрены партитурой и могут вызывать
ложные продвижения HMM. В нашей реализации введён орнаментный
фильтр, основанный на оценке плотности onset’ов: если в окне
$60$\,мс зафиксировано более трёх onset’ов с убывающей энергией,
последовательность считается орнаментом и не передаётся трекеру.
Аналогичная эвристика используется для подавления случайных
повторений нот при двойных нажатиях.

Полный перечень рассмотренных граничных случаев и применяемых правил
сведён в \Tabref{corner-cases}.

\par\addvspace{2pt}
{\centering
\refstepcounter{table}\tablab{corner-cases}%
\footnotesize
\begin{tabular}{@{}p{0.42\linewidth}p{0.5\linewidth}@{}}
  \toprule
  Граничный случай                       & Правило обработки \\
  \midrule
  Фермата / удержание педали             & замораживание перехода, темп до $0{,}7$\,$\nu_0$ \\
  Резкое ритардандо в каденции           & профиль повышенной чувствительности RL \\
  Пропуск $\le 2$ нот                    & разрешённый skip-переход в~\Eqref{transition-matrix} \\
  Пропуск целого такта                   & повторная синхронизация HMM в окне $\pm 4$ ноты \\
  Случайный повтор ноты                  & фильтр двойных нажатий ($<25$\,мс) \\
  Орнамент / мелизм                      & окно $60$\,мс по плотности onset’ов \\
  Фальшивая нота полутоном               & принимается, $b_i(o_t)$ рассчитывается CNN \\
  Октавная ошибка ввода                  & замена кандидата на ближайшую к $p_{\hat\varphi}$ октаву \\
  \bottomrule
\end{tabular}\par
\vspace{2pt}
\tablename~\thetable. Перечень граничных случаев и эвристик их
обработки. Здесь $\nu_0$~--- номинальный темп фрагмента партитуры.\par}
\addvspace{4pt}
